\documentclass[11pt]{article}
\usepackage[T1]{fontenc}
\usepackage{lmodern,microtype,amsmath,amssymb,amsthm,mathtools,mathrsfs}
\usepackage[margin=1in]{geometry}
\usepackage{booktabs,longtable,array,enumitem,tikz-cd}
\usepackage[colorlinks=true,linkcolor=blue,citecolor=blue,urlcolor=blue]{hyperref}
\hypersetup{pdftitle={Paper plan: comparisons, internal theories, and essential incompleteness -- version 1},pdfauthor={Florian Lengyel}}
\newtheorem{theorem}{Theorem}[section]
\newtheorem{lemma}[theorem]{Lemma}
\newtheorem{proposition}[theorem]{Proposition}
\newtheorem{corollary}[theorem]{Corollary}
\theoremstyle{definition}\newtheorem{definition}[theorem]{Definition}
\newcommand{\PA}{\mathrm{PA}}
\newcommand{\N}{\mathbb N}
\newcommand{\Set}{\mathbf{Set}}
\newcommand{\Pfn}{\mathbf{Pfn}}
\newcommand{\Tot}{\operatorname{Tot}}
\newcommand{\Split}{\operatorname{Split}}
\newcommand{\Par}{\operatorname{Par}}
\newcommand{\ran}{\operatorname{ran}}
\newcommand{\Th}{\operatorname{Th}}
\newcommand{\Fm}{\mathsf{Fm}}
\newcommand{\Sent}{\operatorname{Sent}}
\newcommand{\Con}{\operatorname{Con}}
\newcommand{\Prov}{\operatorname{Prov}}
\newcommand{\Tr}{\operatorname{Tr}}
\newcommand{\B}{\mathcal B}
\newcommand{\R}{\mathscr R}
\newcommand{\A}{\mathcal A}
\newcommand{\D}{\mathcal D}
\newcommand{\E}{\mathcal E}
\newcommand{\U}{\mathcal U}
\newcommand{\source}[1]{\par\smallskip\noindent\textit{Source and action.} #1\par\smallskip}
\newcommand{\purpose}[1]{\noindent\textit{Purpose.} #1\par\medskip}
\setlength{\emergencystretch}{3em}
\setlist{topsep=4pt,itemsep=3pt}
\title{Plan for revising\newline\emph{Local connectedness in recursion categories\newline and essential incompleteness}}
\author{Theorem statements, proof outlines, and section dependencies}
\date{Plan version 1 --- 1 November 2026}
\begin{document}
\maketitle

\noindent\textbf{Baseline.} Manuscript version 19, 1 November 2026.\par
\noindent Commit \texttt{2a785a3b5b3b9686e12f72f5118cf14bbbf5e251}.\par
This document plans the next revision. It does not revise the manuscript.
The section and result numbers below are proposed numbers; references
prefixed by ``v19'' identify the current manuscript.

\section*{The proposed argument}
\addcontentsline{toc}{section}{The proposed argument}
The principal theorem will concern a partial conditional in a semantic
target whose morphism lifts to a source recursion category. A Turing
morphism in the source then produces a total morphism whose interpretation
avoids both target domains. Local connectedness and weak choice provide
the target conditional under the existing hypotheses. Range preservation
identifies the interpreted internal theorem object with the image of the
interpreted conclusion morphism.

The arithmetic application will use the fixed ambient category
$S_{\PA}$. Each c.e. theory $A\supseteq\PA$ is represented by an axiom
subobject in $\B_{\PA}=\Tot(\Split S_{\PA})$. Actual consistency of $A$
gives disjoint theorem and refutation domains in the standard
interpretation. A partial program searching for a proof of either a
sentence or its negation supplies the required lift. Neither PA-provable
consistency of $A$ nor PA-provable totality of that search is required.
Recursive separators also give source program lifts; their impossibility
proves recursive inseparability and hence essential incompleteness.

Soundness will have its own theorem. For a represented theory $A$ and
a model $K\models A$, standard numerical proof codes give a conclusion
map factoring through the subobject of sentences valid in $K$.
The numerical code structure and the semantic structure $K$ have
different roles. Taking an ultrapower
of the entire interpretation, including satisfaction, preserves this
factorization for all interpreted proof codes. This includes nonstandard
codes and does not require preservation of initial derivation algebras.
The arithmetic incompleteness theorem for arbitrary consistent extensions
uses consistency directly and therefore does not assume their axioms
true in $\N$.

\paragraph{The consistency claim.}
The soundness argument is a model-theoretic consistency proof: an
independently verified interpretation of the axioms and rules excludes
a proof of contradiction. A formal relative-consistency claim must
specify the theories whose consistency is compared and the metatheory
in which a proof translation or model construction is formalized.
No such formalized argument is supplied here, so this revision will
not claim a new relative-consistency
strength. Ultrapowers transfer a soundness argument; they do not produce
consistency from no semantic assumptions. For arbitrary consistent
extensions in the incompleteness theorem, consistency remains a hypothesis.

\paragraph{The Di Paola--Heller claim.}
Di Paola and Heller explicitly proposed an algebraization of generalized
incompleteness, including inseparability of theorem and refutation sets
\cite{DPH}. The proposed paper supplies a categorical theorem under
specified syntax, representation, interpretation, and lifting hypotheses,
and verifies those hypotheses in concrete recursion categories.
This supports the claim that the paper realizes an inseparability form
of their program under these hypotheses. It does not establish that
every recursion category has internal first-order syntax, that every
internal theory is incomplete, or that all of their proposed generalized
incompleteness statements follow. The arithmetic conclusion is known;
priority for the comparison theorem and the particular realizations
requires comparison with the cited categorical precedents.

\paragraph{Version and scope.}
The syntax appendix keeps its single task: the categorical construction
of syntax and deductively closed theories. Numerical realizations remain
where that appendix is applied. The new soundness and incompleteness
arguments belong in the main text. The three-point application remains
propositional for its selected atom object. The Miller application retains
its distinct claims about uniform computation, real closed fields, and
failure to preserve inseparability.

\begingroup\small
\tableofcontents
\endgroup
\newpage

\section*{Migration from version 19}
\addcontentsline{toc}{section}{Migration from version 19}
\begin{longtable}{@{}>{\raggedright\arraybackslash}p{.17\textwidth}>{\raggedright\arraybackslash}p{.23\textwidth}>{\raggedright\arraybackslash}p{.53\textwidth}@{}}
\toprule
Current location & Proposed location & Mathematical action\\\midrule\endhead
\S\S1--2 & \S\S1--2 & Rewrite the introduction after the results are settled; retain categorical conventions and definitions.\\
\S3 & \S3 and \S4 & Prove the comparison theorem first for partial target conditionals. Retain local connectedness as a sufficient condition and the source inseparability corollaries.\\
\S4, results 4.1--4.6 & \S8 & Retain the Miller construction, reduction, syntax realization, and theory theorem. Change only the route through the general results.\\
\S5, results 5.1--5.9 & \S5 & Retain the realization in $\B_T$ and all nine supporting statements. Use distinct letters for the ambient theory $B$ and represented theory $A$.\\
\S6, result 6.1 & \S6 & Retain one-model interpretation; add standard and ultrapower soundness.\\
\S7, results 7.1--7.3 & \S6, \S7, Appendix D & Keep one-model source inseparability in \S6. Generalize the search lemma for \S7. Move the faithful family and theory-extension square to Appendix D.\\
\S8, results 8.1--8.2 & \S5 and \S7 & Retain the $\mathrm{Rec}$ realization as a comparison. Make the arithmetic derivation from $S_{\PA}$ the main proof.\\
\S9, results 9.1--9.6 & \S9 & Retain propositional initiality, endpoint computations, consistent theories, and endpoint essential incompleteness.\\
Appendix A\newline A.1--A.10 & Appendix B & Retain the complete three-point cosheaf construction and its recursion-category verification.\\
Appendix B\newline B.1--B.7 & Appendix A & Retain the categorical syntax and theory construction, its substitution dependencies, and the range/image lemma.\\
Appendix C, C.1 & Appendix C & Retain the counterexample to dropping lifting; identify exactly which proposed corollary it refutes without lifting.\\
\bottomrule
\end{longtable}

\section{Introduction}
\purpose{State the comparison theorem, its logical application, and the verified examples before discussing its relation to the Di Paola--Heller program.}
The first mathematical claim should describe the total morphism avoiding
two disjoint interpreted domains when their reversed partial conditional
lifts. The following paragraph should explain that internally generated
theories provide the theorem domains, while a semantic interpretation
provides consistency. Then state the arithmetic, Miller, and cosheaf
conclusions with their actual computability hypotheses.

Keep the distinction between this theorem and Heller's existence theorem:
the source here already has a Turing morphism. State that local
connectedness is a sufficient target hypothesis. The target need not be
a recursion category, and the functor need not be faithful or preserve
restriction products. Describe the faithful family of models only as a
separate representation result.

The related-work discussion should place Di Paola--Heller, Bernardi,
Montagna, and Di Paola--Montagna before the present construction
\cite{DPH,Bernardi,Montagna,DPM}. The categorical syntax precedents remain
in Appendix A with their precise functions. Miller's construction and the
cosheaf literature stay with their applications. No claim of originality
should rest on the fact that a theorem has been reformulated in the
present notation.
\source{Replace v19 \S1 and its roadmap after \S\S3--7 are revised. Draft the abstract last. The current title can be retained because local connectedness still supplies a general sufficient condition.}

\section{Restriction categories, ranges, and recursion}
\purpose{Fix the structures and conventions used by the comparison theorem and by the theory construction.}
Retain the definitions in v19 \S2: restriction idempotents and total
morphisms; restriction zeros; ranges and their agreement property;
restriction products and coproducts; inverse images; joins and
complementarity; partial monomorphisms, sections, and weak choice;
local connectedness; restriction splitting; recursion categories and
Turing morphisms. Write composition as $gf=g\circ f$.

For $f:X\rightharpoonup Y$ and a restriction idempotent $e$ on $Y$,
$f^*e=\overline{ef}$. A range functor preserves restrictions, ranges,
restriction zeros, and binary restriction coproducts. It consequently
preserves joins because $a\vee b=\ran[a,b]$. Write
\[
 \chi_{X,Y}:FX+FY\xrightarrow{\cong}F(X+Y)
\]
for its coproduct comparison. The definition does not include preservation
of restriction products.

For an original source object $U$, the inclusion
$\R\to\Split\R$ is full on the hom-sets between original objects.
Thus a morphism $U\rightharpoonup U+U$ constructed in the splitting is
also a source morphism. Formula, sentence, and proof objects may belong
to $\Tot(\Split\R)$; this must be stated whenever they are introduced.
\source{Retain v19 \S2. Add only the conventions needed to distinguish $F$ from its extension to restriction splittings and its restriction to total categories. Use Cockett--Lack and Cockett--Hofstra for the established categorical framework \cite{CockettLack,CockettHofstra}.}

\section{Diagonal arguments and lifted partial conditionals}
\purpose{Prove the comparison theorem before imposing local connectedness.}

\begin{theorem}[Comparison theorem]\label{thm:comparison}
Let $\R$ be a recursion category with ranges and binary restriction
coproducts, and let $\A$ be a restriction category with ranges,
restriction zeros, and binary restriction coproducts. Let
$F:\R\to\A$ be a range functor. Choose a Turing morphism
$\tau:U\times U\rightharpoonup U$ and an isomorphism
$b:U\to U+U$. Put
\[
 h=b\tau\Delta_U,\qquad
 a_0=h^*(1_U+0_U),\qquad a_1=h^*(0_U+1_U).
\]
Let $P,N$ be disjoint restriction idempotents on $FU$ such that
$F(a_0)\le P$ and $F(a_1)\le N$. Suppose a partial morphism
$s:FU\rightharpoonup FU+FU$ satisfies
\[
 s^*(1_{FU}+0_{FU})=N,\qquad
 s^*(0_{FU}+1_{FU})=P,
\]
and suppose there is $\ell:U\rightharpoonup U+U$ with
$F\ell=\chi_{U,U}s$.
Then there is a total $g:U\to U$ such that
\[
 (Fg)^*P=(Fg)^*N=0_{FU}.
\]
If $F(1_U)\ne0$ and inverse image along $Fg$ preserves the join
$P\vee N$, then $P\vee N\ne1_{FU}$.
\end{theorem}
\begin{proof}
Put $u=b^{-1}\ell:U\rightharpoonup U$. Apply the Turing property
to $u\pi_2:U\times U\rightharpoonup U$, obtaining a total
$g:U\to U$ with $\tau(g\times1_U)=u\pi_2$.
Precompose with $\langle1_U,g\rangle:U\to U\times U$ to obtain
$\tau\Delta_Ug=ug$, hence $hg=\ell g$.

Write $E=Fg$ and $H=\chi_{U,U}^{-1}Fh$. Then $HE=sE$.
Preservation of restrictions and the coproduct comparison give
\[
 E^*F(a_0)=E^*N,\qquad E^*F(a_1)=E^*P.
\]
The two inequalities in the statement give $E^*N\le E^*P$
and $E^*P\le E^*N$. Their meet is
$E^*(PN)=0$, so both are zero. Totality of $E$ and preservation
of the indicated join would give
$1_{FU}=E^*1_{FU}=E^*(P\vee N)=0$ if $P\vee N=1_{FU}$.
\end{proof}
\source{New statement extending the proof of v19 Theorem 3.3. The partial conditional need not be total, its source lift need not be total, and neither faithfulness nor product preservation is used. Range preservation is retained as the paper's common comparison structure, although the displayed diagonal calculation only uses restrictions, zeros, and coproducts.}

\begin{corollary}[Separators whose conditionals lift]\label{cor:separator}
Under the categorical assumptions of Theorem~\ref{thm:comparison},
assume $F(1_U)\ne0$ and that total inverse images in $\A$ preserve
binary joins. A complementary pair $P,N$ with
$F(a_0)\le P$ and $F(a_1)\le N$ cannot have a reversed conditional
admitting a source lift as in that theorem.
\end{corollary}
\begin{proof}
Apply the theorem and use $P\vee N=1_{FU}$.
\end{proof}
This does not assert inseparability among all target restriction
idempotents. In $\Pfn$ every disjoint pair has a set-theoretic
separator. The obstruction concerns separators whose conditional has
the specified source lift.

\begin{lemma}[Connected domains and coproducts]\label{lem:connected}
Suppose $\A$ has the restriction structure above and inverse images
preserve binary joins. If $a$ is a connected restriction idempotent
on $X$ and $f:X\to Y+Z$ is total, then
$(1_Y+0_Z)fa=fa$ or $(0_Y+1_Z)fa=fa$.
\end{lemma}
\begin{proof}
Pull the disjoint summand decomposition back along $f$ and meet
with $a$. Connectedness makes one piece zero; the other is $a$.
The restriction identity gives the asserted equation.
\end{proof}
\source{Retain v19 Lemma 3.1.}

\begin{proposition}[Conditionals in a locally connected category]\label{prop:conditionals}
Let $\A$ be split and locally connected, with ranges, restriction
zeros, and binary restriction coproducts. Assume weak choice and
preservation of binary joins by inverse images. For disjoint
restriction idempotents $P,N$ on $X$, the partial morphism
\[
 p=\nabla_X(N+P):X+X\rightharpoonup X
\]
has a unique section $s:X\rightharpoonup X+X$. It satisfies
\[
 \bar s=P\vee N,\qquad
 s^*(1_X+0_X)=N,\qquad s^*(0_X+1_X)=P.
\]
\end{proposition}
\begin{proof}
Use Lemma~\ref{lem:connected} on connected domain restrictions to
show that $p$ is a partial monomorphism. Split the common domain
of $pu,pv$ to extend the argument from total maps to arbitrary
partial maps. Weak choice supplies a section. The section equations
give $\bar s=\ran(p)=P\vee N$; partial monicity gives uniqueness.
The identities $sP=i_1P$ and $sN=i_0N$, together with this domain
equation, give the two summand inverse images.
\end{proof}
\source{Retain v19 Proposition 3.2 and state its partial summand equations explicitly. This proposition is where local connectedness and weak choice are used.}

\begin{corollary}[Inseparability in the source]\label{cor:source}
Let $F:\R\to\A$, $\tau$, $b$, and $a_0,a_1$ be as in
Theorem~\ref{thm:comparison}, with $\A$ as in
Proposition~\ref{prop:conditionals} and $F(1_U)\ne0$.
If the reversed conditional associated with $Fc,Fc'$ lifts for every
complementary source pair $c,c'$ on $U$, then $a_0,a_1$ are inseparable
in $\R$.
\end{corollary}
\begin{proof}
A source separator is sent to a complementary target pair, since
a range functor preserves joins, meets, and zero. The assumed lift
contradicts Corollary~\ref{cor:separator}.
\end{proof}
\source{Retain the conclusion of v19 Theorem 3.3 as a corollary. Appendix C shows that its lifting hypothesis cannot simply be omitted.}

\begin{proposition}[Disjoint joins give source lifts]\label{prop:joins}
If parallel partial morphisms with disjoint domains have joins in
the source, then for complementary $c,c'$ on $X$ the morphism
\[
 \ell=(i_1c)\vee(i_0c'):X\to X+X
\]
is total and satisfies $\ell c=i_1c$ and $\ell c'=i_0c'$.
If $F:\R\to\A$ is a range functor and $\A$ satisfies
Proposition~\ref{prop:conditionals}, then
$\chi_{X,X}^{-1}F\ell$ is the target section of that proposition.
\end{proposition}
\begin{proof}
The join equations give the restrictions on $c,c'$; their union
is the full domain. After the coproduct identification, both maps agree on
$Fc,Fc'$. Their copairing has full range, so range agreement gives
equality.
\end{proof}
\source{Retain v19 Proposition 3.4. It remains useful in the Miller and numerical applications.}

\section{Internal theories and semantic consistency}
\purpose{Connect the abstract construction in Appendix A to the comparison theorem, without identifying syntax with an arbitrary restriction idempotent.}
Let $\D=\Tot(\Split\R)$ satisfy the hypotheses used in Appendix A.
Assume the semantic target $\A$ is split, and choose an extension
$\widehat F:\Split\R\to\A$ preserving ranges and the restriction
structure. Use the canonical functor to $\Split\A$ followed by
chosen splittings in $\A$, with the identity splitting on original
objects, so that $\widehat F|_{\R}=F$.
Let $\Phi\in\D$ be the formula object and $j:S\hookrightarrow\Phi$
its sentence subobject. For an axiom subobject $A\hookrightarrow S$,
let $c_A:P_A\to\Phi$ be the initial derivation algebra's conclusion.
Form the pullback in $\D$:
\[
\begin{tikzcd}
 P_{A,S}\arrow[r]\arrow[d,"c_{A,S}"']&P_A\arrow[d,"c_A"]\\
 S\arrow[r,hook,"j"']&\Phi.
\end{tikzcd}
\]
Put $t_A=\ran(c_{A,S})$ in $\Split\R$ and let
$\nu:S\to S$ be syntactic negation. These objects and maps are
defined by the appendix; the application must prove their realization.

\begin{proposition}[Interpretation of the theorem object]\label{prop:image}
Let $F:\R\to\A$ and $\widehat F$ be as specified above. Then
\[
 \widehat F(t_A)=\ran(\widehat F(c_{A,S})).
\]
If the target is presented as $\Par(\E,\operatorname{Mono})$ for
a regular category $\E$, a factorization in $\E$ of the total
conclusion map through a validity monomorphism $V\hookrightarrow
\widehat F S$ implies that the image of the internal theorem object
lies in $V$.
\end{proposition}
\begin{proof}
The first equation is preservation of ranges. Apply the range/image
lemma of Appendix A to the conclusion map in the total category;
factorization through a monomorphism bounds its regular image.
\end{proof}
No claim that $\widehat F(c_A)$ remains an initial derivation algebra
is required. The image of $\widehat F(c_{A,S})$ is the interpreted
sentence-theorem object, which may contain conclusions of nonstandard
proof codes.

\begin{theorem}[Incompleteness from semantic consistency]\label{thm:logical}
Use the source data of Theorem~\ref{thm:comparison} and suppose the
target has total inverse images preserving binary joins, with
$F(1_U)\ne0$. Let $\theta:U\to S$ be total in $\Split\R$.
Put $t=\widehat F(t_A)$ and define restriction idempotents on $FU$ by
\[
 P=(\widehat F\theta)^*t,\qquad
 N=(\widehat F\theta)^*(\widehat F\nu)^*t.
\]
Assume $t(\widehat F\nu)^*t=0$, $F(a_0)\le P$, and
$F(a_1)\le N$. If the reversed partial conditional for $P,N$
has a source lift, there is a total $g:U\to U$ such that
\[
 (\widehat F(\theta g))^*t=0,\qquad
 (\widehat F(\theta g))^*(\widehat F\nu)^*t=0.
\]
Consequently $t\vee(\widehat F\nu)^*t\ne1_{\widehat F S}$.
\end{theorem}
\begin{proof}
Target consistency makes $P,N$ disjoint. Apply
Theorem~\ref{thm:comparison}. The two displayed equations are
its conclusion after composing inverse images. If the two theorem
domains covered $\widehat F S$, their pullbacks along the total
$\widehat F(\theta g)$ would cover $FU$, contradicting
$F(1_U)\ne0$.
\end{proof}
\source{New logical application of the partial comparison theorem. It uses consistency in the target and does not require a source consistency equation. The two representation inequalities are target inequalities; an application may establish them directly or obtain them by interpreting source inequalities.}

\begin{corollary}[The source-consistency criterion]\label{cor:source-logical}
Let a restriction category have zeros and binary joins preserved by
total inverse images. Let $a,b$ be inseparable on $U$, let
$\theta:U\to S$ and $\nu:S\to S$ be total, and let $t_0$ be a
restriction idempotent on $S$ satisfying
$a\le\theta^*t_0$ and $b\le\theta^*\nu^*t_0$.
Every restriction idempotent $t$ on $S$ with $t\ge t_0$
and $t\nu^*t=0$ has inseparable theorem and refutation domains
$t,\nu^*t$; in particular, $t\vee\nu^*t\ne1_S$.
\end{corollary}
\begin{proof}
A complementary separator of $t,\nu^*t$ pulls back to one of
$a,b$. The same argument applies when $t,\nu^*t$ themselves are
assumed complementary.
\end{proof}
\source{Retain v19 Corollary 3.5 for the Miller and cosheaf arguments, where source consistency is actually established. Do not use it for an arithmetic theory's own proof object in $S_T$ merely from external consistency.}

\section{Syntactic recursion categories and their internal syntax}
\purpose{Prove the appendix's realization where it is used.}
Throughout this section, $B\supseteq\PA$ denotes the ambient
arithmetic theory and $A$ denotes a represented axiom presentation.
This replaces the ambiguous reuse of $T$ for both roles.

\begin{lemma}[Arithmetic computation laws]\label{lem:computation}
Fix the program enumeration and numerical product and coproduct
codings of v19 \S5. The graph predicates are $\Sigma_1$ and
PA-provably functional. PA proves the uniform graph identities for
composition, identities, projections, pairing, coproduct injections,
copairing, domains, and the chosen universal evaluator. In particular,
there are a universal index $u$ and primitive-recursive functions
$\operatorname{s},\operatorname{comp}:\N\times\N\to\N$ for which
PA proves the universal closures of
\[
\begin{aligned}
G_u(\langle e,x\rangle,y)&\leftrightarrow G(e,x,y),\\
G(\operatorname{s}(e,a),x,y)&\leftrightarrow G(e,\langle a,x\rangle,y),\\
G(\operatorname{comp}(f,e),x,z)&\leftrightarrow
 \exists y\,[G(e,x,y)\land G(f,y,z)].
\end{aligned}
\]
For each fixed $e$, the function $a\mapsto\operatorname{s}(e,a)$
has a program whose graph and totality PA proves.
\end{lemma}
\begin{proof}
Retain the formalized coding argument used in v19 Lemma 5.1,
including its dependence on Montagna's computation laws
\cite{Montagna}, and verify the transfer to the chosen trace predicates.
Composition concatenates two traces and records the intermediate
value. The structural programs implement their numerical codings.
The evaluator simulates the indexed program, and the parameter
program gives its uniform factorization. All equalities are universal
graph equalities, not separate proofs for individual numerals.
\end{proof}
\source{Retain v19 Lemma 5.1 and the definitions surrounding it. Morphisms of $S_B$ are indices modulo $B$-provable uniform graph equality. The primed pointwise quotients are not substituted for this category.}

\begin{lemma}[Functional search]\label{lem:search}
For a $\Sigma_1$ relation $\Gamma(x,y)$ there is a partial recursive
search program with domain $\exists y\,\Gamma(x,y)$; if $B$ proves
functionality of $\Gamma$, it proves uniform equality of that
program's graph with $\Gamma$.
\end{lemma}
\begin{proof}
Search increasing finite bounds for a witness to the primitive
recursive matrix. The trace verification proves existence exactly
on the existential domain. Under functionality, any discovered
witness has the unique permitted output.
\end{proof}
\source{Retain v19 Lemma 5.2.}

\begin{lemma}[The total splitting]\label{lem:total}
For consistent $B\supseteq\PA$, the category
$\B_B=\Tot(\Split S_B)$ is equivalent to the category of
$\Sigma_1$ carrier predicates and $B$-provably total functional
$\Sigma_1$ relations, modulo $B$-provable uniform graph equality.
It is regular with disjoint pullback-stable finite coproducts.
Every monomorphism is an admissible domain monomorphism, and
\[
 \Par(\B_B,\operatorname{Mono})\simeq\Split S_B.
\]
\end{lemma}
\begin{proof}
Use Lemma~\ref{lem:search} to identify functional relations with
program morphisms. Products and equalizers use conjunction and
numerical pairing; tagged unions give extensive coproducts. Images
are existential graph predicates. A graph whose surjectivity is proved
in $B$ has a witness-search section, so it is the coequalizer of its
kernel pair and remains so
under pullback. For a monomorphism, its kernel-pair projections are
equal; this proves injectivity uniformly in $B$. Searching its
transposed graph gives its partial inverse. These constructions give
the asserted partial-map equivalence.
\end{proof}
\source{Retain all parts of v19 Lemma 5.3, including the categorical proof of injectivity; do not replace it by a standard-points argument.}

\begin{lemma}[Finite membership conditions]\label{lem:finite}
PA proves that a coded finite family of $\Sigma_1$ membership
conditions holds exactly when a coded finite family of witnesses
exists. Consequently the finite validity tests used below are
$\Sigma_1$ predicates with the required PA-provable equivalences.
\end{lemma}
\begin{proof}
Induction on the coded list length appends a witness at each step.
Conversely each component of a witness list verifies its membership
condition. This is induction in PA, so it covers internally finite
lists in every model of $B$.
\end{proof}
\source{Retain v19 Lemma 5.4.}

\begin{proposition}[Syntax and lists with parameters]\label{prop:syntax}
Fix bounded finite arity and variable and signature objects of
$\B_B$. The numerical term and formula tree objects realize the
appendix's initial syntax algebra, stable under every parameter
object. The numerical list object realizes the corresponding
parameterized initial list algebra.
\end{proposition}
\begin{proof}
Use canonical trees and lists as carriers. Numerical equality makes
every represented variable carrier decidable in $\B_B$.
For an arbitrary algebra
over a parameter object, an evaluation trace records the intermediate
values, the algebra-operation witnesses, and the common parameter.
Lemma~\ref{lem:finite} makes validity $\Sigma_1$. Induction in $B$
proves existence and uniqueness of evaluation, and a second induction
proves uniform graph equality with every algebra morphism. This gives
initiality in each slice. It is not merely the initiality of ordinary
finite trees in $\Set$.
\end{proof}
\source{Retain v19 Propositions 5.5 and 5.6 as separate statements in the manuscript; grouped here because their proof mechanism is shared.}

\begin{proposition}[Substitution realization]\label{prop:substitution}
The numerical free-variable, free-for-substitution, term-substitution,
and formula-substitution operations realize the categorical folds and
equations specified in Appendix A in $\B_B$.
\end{proposition}
\begin{proof}
Verify the constructor equations, including the quantifier case and
the free-for condition, in PA. Parameterized initiality identifies
these operations with the appendix's folds; uniqueness gives the
claimed identities as $B$-provable graph equalities.
\end{proof}
\source{Retain v19 Proposition 5.7. This proof stays in the realization section, not Appendix A.}

\begin{lemma}[Initial derivations]\label{lem:proof-initial}
Every bounded-arity rule presentation by objects and total morphisms
of $\B_B$ has an initial derivation algebra in the slice over its
formula object.
\end{lemma}
\begin{proof}
Use trees labeled by actual rule instances with ordered children.
Finite membership conditions and the matching of premises to child
conclusions are $\Sigma_1$. Evaluation into any rule algebra is
$\Sigma_1$; $B$-induction proves its existence, functionality,
algebra equations, and uniqueness. Membership witnesses verify labels;
they are not extra constructors of the derivation object.
\end{proof}
\source{Retain v19 Lemma 5.8.}

\begin{theorem}[Internally generated theories]\label{thm:internal}
For consistent $B\supseteq\PA$, the bounded-arity first-order
structure of Appendix A is realized in $\B_B$. Every axiom
subobject $A\hookrightarrow S$ has an initial derivation algebra.
Its conclusion image is the least subobject closed under the
specified rules, and its sentence pullback is the generated sentence
theory. Let $e_A$ be the program whose domain is exactly the
valid-derivation predicate $P_A$ and whose output is the conclusion.
The partial identity on the numerical codes of the sentence theory,
extended by zero outside the sentence carrier, is
\[
 t_A^{\omega}=[\Sent(z)\land\exists p\,G_{e_A}(p,z)]_B:
 \omega\rightharpoonup\omega.
\]
Its restriction to the sentence splitting is
$t_A:S\rightharpoonup S$.
\end{theorem}
\begin{proof}
Apply Lemmas~\ref{lem:total} and \ref{lem:proof-initial},
Propositions~\ref{prop:syntax} and \ref{prop:substitution}, and the
least-theory theorem of Appendix A. Functional search realizes the
sentence-conclusion pullback. The range/image lemma identifies its
program range with the regular image in the total category.
\end{proof}
\source{Retain v19 Theorem 5.9. Explicitly instantiate it with ambient $B=\PA$ and an independently chosen c.e. axiom presentation $A$, using the carrier predicate $\Sent(z)\land\operatorname{Ax}_A(z)$. The sentence conjunct makes this a typed subobject even if the enumeration's output typing has not been proved in PA.}

\begin{proposition}[Numerical realization]\label{prop:rec-realization}
The category $\Tot(\Split\mathrm{Rec})$ is equivalent to the category
of c.e. numerical carriers and partial recursive functions total on
their source carriers. It realizes the same bounded-arity first-order
syntax, substitution, lists, and generated-theory construction.
\end{proposition}
\begin{proof}
Finite limits use pairing and equality, coproducts use tags, and
images are c.e. by dovetailing. Evaluation of ordinary finite trees
gives initiality and uniqueness by induction. Finite nonlogical
axiom membership tests are semidecidable. The range/image lemma gives
the theorem partial identity.
\end{proof}
\source{Retain v19 Proposition 8.1 as the numerical comparison. The main arithmetic proof below will additionally use the internal theory object in $\B_{\PA}$.}

\paragraph{Computational inseparability and completeness.}
For $B=\Th(\N)$, interpreting each program in $\N$ gives an
isomorphism $S_{\Th(\N)}\cong\mathrm{Rec}$ with the chosen common
object presentation. Two indices are equal in the source exactly
when their uniform graph-equivalence sentence belongs to
$\Th(\N)$, which is equivalent to equality of their standard
partial functions. Thus $S_{\Th(\N)}$ has inseparable computational
pairs although $\Th(\N)$ is complete. This example distinguishes
the existence of an inseparable pair from its representation by
the theorem and refutation domains of a specified internal theory.
\source{Retain the unnumbered example from v19 \S8.}

\section{Model interpretation, soundness, and ultrapowers}
\purpose{Construct the semantic comparison and prove soundness for its entire interpreted internal proof object.}

\begin{proposition}[Interpretation in one model]\label{prop:model}
Let $B\supseteq\PA$ be consistent and let $M\models B$.
Interpreting the program graph formulas defines a range functor
$F_M:S_B\to\Pfn$ preserving restriction products. Its extension
to splittings, with the subset splittings in $\Pfn$, restricts to
\[
 G_M:\B_B\longrightarrow\Set
\]
preserving finite limits, finite coproducts, and regular images.
\end{proposition}
\begin{proof}
Interpret the uniform computation identities of
Lemma~\ref{lem:computation} in $M$. $B$-provable graph equality
makes the assignment independent of indices. In the predicate
presentation, a carrier $C$ becomes
$C(M)=\{x\in M:M\models C(x)\}$. The finite-limit formulas,
tagged coproducts, and existential image formulas become the
corresponding set constructions. Internal traces and witnesses range
over all of $M$, including its nonstandard elements.
\end{proof}
\source{Retain v19 Proposition 6.1 and add the total-category preservation statement with this proof. A single model suffices here; no faithful family is needed.}

\begin{proposition}[Source complementarity and inseparability]\label{prop:st-pair}
For complementary restriction idempotents $c,c'$ on $X$ in $S_B$,
there is a total $s:X\to X+X$ with $sc=i_1c$ and $sc'=i_0c'$.
For any $M\models B$, its image is the reversed conditional for
$F_Mc,F_Mc'$. Consequently the diagonal pair is inseparable in
$S_B$ for every consistent $B\supseteq\PA$.
\end{proposition}
\begin{proof}
Search for a witness to either $\Sigma_1$ domain predicate.
$B$-proved disjointness and coverage give the uniform graph identity
and totality in the source. Interpretation gives the required
conditional. The target $\Pfn$ satisfies
Proposition~\ref{prop:conditionals}, and $F_M(1_U)\ne0$.
Apply Corollary~\ref{cor:source}. A model exists by completeness.
\end{proof}
\source{Retain the substantive conclusions of v19 Lemma 7.2 and Corollary 7.3. Faithfulness of the all-model product is not a dependency. Theory-extension naturality is recorded in Appendix D.}

\begin{proposition}[Soundness from one semantic model]\label{prop:soundness}
Fix the ambient category $\B_{\PA}$, a c.e. first-order arithmetic
axiom presentation $A$, its sentence-conclusion map
$c_{A,S}:P_{A,S}\to S$, and a model $K\models A$.
Let $\nu:S\to S$ be syntactic negation. For a standard sentence
code $s$, let $\varphi_s$ denote the sentence it codes.
Define a subobject in $\Set$ by
\[
 V_K=\{s\in G_{\N}S:K\models\varphi_s\},\qquad
 v_K:V_K\hookrightarrow G_{\N}S.
\]
There is a function $e_K:G_{\N}P_{A,S}\to V_K$ with
$v_Ke_K=G_{\N}c_{A,S}$. Moreover
$V_K\cap(G_{\N}\nu)^{-1}V_K=\varnothing$, and $A$ is consistent.
\end{proposition}
\begin{proof}
In the standard numerical interpretation, every proof code represents
an actual finite derivation in the fixed calculus. Its nonlogical
axioms hold in $K$; the logical axiom schemes are valid; modus ponens
and generalization preserve validity. Induct on the entire derivation,
taking validity of an open formula to mean truth under every
assignment in $K$. At a sentence conclusion this is
$K\models\varphi_s$, which gives the factorization.
A sentence and its negation cannot both
hold in $K$. Since every actual proof occurs in this carrier, a
pair of contradictory proofs would contradict that disjointness.
\end{proof}
\source{New main-text statement. This is ordinary soundness expressed as a factorization of the interpreted internal proof map. $K$ is a semantic structure for $A$, whereas $G_{\N}$ interprets numerical codes. It is not $G_K$. If $K$ is constructed from assumed consistency by completeness, this is not an independent proof of that consistency.}

\begin{lemma}[The ultrapower functor]\label{lem:ultrafunctor}
Work in a set-theoretic metatheory with choice. Let $\U$ be a proper
ultrafilter on a nonempty set $I$. For a set $X$ and two families
$(x_i),(y_i)\in X^I$, define
\[
 (x_i)\sim_{\U}(y_i)
 \quad\Longleftrightarrow\quad
 \{i\in I:x_i=y_i\}\in\U.
\]
The assignments
\[
 \mathscr U(X)=X^I/{\sim_{\U}},\qquad
 \mathscr U(f)([x_i])=[f(x_i)]
\]
for sets $X$ and functions $f:X\to Y$ define a faithful functor
$\mathscr U:\Set\to\Set$ preserving
finite limits, finite coproducts, surjections, and regular images.
Here $X^I$ is the set of functions $I\to X$ in $\Set$.
\end{lemma}
\begin{proof}
Composition is coordinatewise; constant families detect inequality
of functions. The empty set and a singleton are fixed by $\mathscr U$.
Pairing proves finite-product preservation. An
equalizer condition holds on a $\U$-large set; replace the other
coordinates by a fixed element of that equalizer to obtain a
representative in its ultrapower. If the equalizer is empty, both
sets are empty. For a binary coproduct, the ultrafilter selects one
of the two tag sets, and the same replacement gives a representative
entirely in that summand. Together with the empty-coproduct case,
this proves finite-coproduct preservation.
Surjections have sections by choice, and their
section equations are preserved. Surjection--injection factorizations
therefore remain image factorizations.
\end{proof}
\source{New supporting lemma. Cite Keisler's functorial ultrapower construction and prove the finite-limit and finite-coproduct assertions used here \cite[Theorem 3.1 and Proposition 4.2]{Keisler}. No exponentials or $I$-indexed powers are asserted in $\B_{\PA}$.}

\begin{theorem}[Ultrapower comparison]\label{thm:ultra}
Let $\U$ be a proper ultrafilter on a nonempty set $I$, and let
$\mathscr U$ be the functor of Lemma~\ref{lem:ultrafunctor}.
Let $B\supseteq\PA$ satisfy $\N\models B$, and put
$M=\N^I/\U$. There is a natural isomorphism
\[
 \eta:\mathscr U G_{\N}\xrightarrow{\cong}G_M
 \quad\text{on }\B_B.
\]
It identifies the ultrapowers of the standard sentence-conclusion
maps and their images with their interpretations in $M$.
\end{theorem}
\begin{proof}
For a carrier predicate $C$, \L{}o\'s's theorem identifies
$C(M)$ with classes having $C(\N)$-valued coordinates on a large
set. Replace other coordinates by a fixed member of $C(\N)$;
the empty case gives empty sets on both sides. For a functional
graph $\Gamma$, \L{}o\'s identifies its interpreted graph with the
ultrapower graph. This proves naturality. Finite-limit and image
preservation give the conclusion about sentence pullbacks and
theorem images \cite{Keisler}.
\end{proof}
\source{New statement. Initial-algebra preservation is not asserted: for $I=\N$ and $\U$ nonprincipal, the ultrapower of $(\N,0,\mathrm{succ})$ is not the initial $1+X$-algebra in $\Set$.}

\begin{theorem}[Soundness for the full interpreted proof object]\label{thm:ultrasound}
Under Proposition~\ref{prop:soundness}, let $I,\U,\mathscr U$
be as in Lemma~\ref{lem:ultrafunctor}, put $M=\N^I/\U$,
and use Theorem~\ref{thm:ultra} with ambient $B=\PA$.
There is a factorization
\[
 G_MP_{A,S}\xrightarrow{e_{K,\U}}\mathscr U(V_K)
 \xrightarrow{v_{K,\U}}G_MS,
 \qquad v_{K,\U}e_{K,\U}=G_Mc_{A,S}.
\]
Its validity subobject and its inverse image under $G_M\nu$ are
disjoint. Hence the interpreted internal theorem and refutation
domains are disjoint, including their nonstandard codes.
\end{theorem}
\begin{proof}
Apply $\mathscr U$ to $v_Ke_K=G_{\N}c_{A,S}$ and use the
components of $\eta$. Explicitly,
\[
 e_{K,\U}=\mathscr U(e_K)\eta_{P_{A,S}}^{-1},\qquad
 v_{K,\U}=\eta_S\mathscr U(v_K).
\]
Naturality gives the equation, and preservation of pullbacks and
the empty set gives disjointness. Proposition~\ref{prop:image}
identifies the interpreted internal theorem object with this
conclusion image. No induction over externally finite target
derivations is being substituted for this factorization.
\end{proof}
\source{New theorem. For full first-order satisfaction, ultrapower the many-sorted structure containing standard numerical codes, the semantic domain $K$, assignments, and satisfaction. The code sort becomes $\N^I/\U$ and the semantic sort becomes $K^I/\U$. Their roles must remain distinct. The special case $K=\N\models A$ gives the expanded-standard-model construction discussed in the continuation.}

\paragraph{The assignment sort.}
One possible standard assignment sort is $K^{\N}$ in $\Set$.
Its transferred sort is $\mathscr U(K^{\N})$; it is not identified
with $(\mathscr U K)^{\mathscr U\N}$. Include evaluation, assignment
modification, and the satisfaction relation before taking the
ultrapower. Their compositional clauses for connectives,
quantification, and substitution transfer by \L{}o\'s's theorem.
The predicate obtained is part of the semantic expansion, not an
arithmetically definable truth predicate in the source.

\begin{corollary}[Transfer of an assumed consistency]\label{cor:consistency}
Let $I,\U,\mathscr U$ be as in Lemma~\ref{lem:ultrafunctor}.
Let $A$ be any consistent c.e. arithmetic theory represented in
$\B_{\PA}$, and put $M=\N^I/\U$. Its interpreted theorem and
refutation domains are disjoint. No assumption $\N\models A$
is required.
\end{corollary}
\begin{proof}
The standard conclusion image is the actual theorem set of $A$,
so its negation intersection is empty by consistency. Apply the
finite-limit and image preservation of $\mathscr U$ and the
comparison isomorphism. Equivalently, transfer the true arithmetic
sentence asserting that no sentence has both an $A$-proof and an
$A$-refutation. A semantic model is not needed for this shorter proof.
\end{proof}
\source{New corollary used in the arithmetic incompleteness theorem. This transfers a hypothesis; it is not a proof of its premise.}

\begin{proposition}[Limits of model interpretation]\label{prop:limits}
The section-lifting hypotheses do not imply soundness for every model
interpretation. Moreover, taking an ultrapower of a fixed model does
not improve faithfulness of its graph interpretation.
\end{proposition}
\begin{proof}
Assuming consistency of PA, choose
$L\models\PA+\neg\Con(\PA)$, with the proof predicate for the
chosen calculus. The interpretation $G_L$ contains a proof code whose
conclusion is $0=1$. Its conclusion map cannot factor through any
validity subobject excluding that sentence, although the source
complementary-conditionals lift as in Proposition~\ref{prop:st-pair}.
Every ultrapower of $L$ retains this sentence. The example is also
used in the partial-combinatory-algebra literature \cite[\S7]{Terwijn}.

For the faithfulness assertion, equality of two interpreted program
graphs is expressed by their uniform graph-equality sentence. A
model and its ultrapower satisfy the same such sentences. Thus
their functors identify exactly the same pairs of source morphisms.
\end{proof}
\source{Retain this as a short limitation proposition adjacent to soundness. It prevents substituting an arbitrary $G_L$ for the comparison of Theorem~\ref{thm:ultrasound}. It does not obstruct that theorem.}

\section{Arithmetic: theorem sets, refutations, and essential incompleteness}
\purpose{Apply the categorical comparison theorem in the fixed ambient category $S_{\PA}$ to every consistent c.e. extension of PA.}
Fix the arithmetic signature, the internal syntax of
Theorem~\ref{thm:internal}, and its sentence object $S\in\B_{\PA}$.
For every c.e. axiom presentation $A$, the standard interpretation
of its generated-theory subobject is the set of actual finite
consequences of $A$.
This identification follows from the exact carrier predicate for
derivation trees, not merely from an abstract claim of initiality.
Write $\Prov_A(z)$ for the $\Sigma_1$ sentence-conclusion image
predicate supplied by Theorem~\ref{thm:internal}, including the
sentence predicate and the chosen axiom presentation.

\begin{lemma}[Representation of the diagonal domains]\label{lem:reduction}
Choose the numerical diagonal domains $D_0,D_1$ obtained by
interpreting $a_0,a_1$ in $\N$. There are $\Sigma_1$ formulas
$\alpha(x),\beta(x)$ and a total morphism in $\Split S_{\PA}$
\[
 \theta:\omega\to S,\qquad
 (\widehat F_{\N}\theta)(n)=\ulcorner\alpha(\bar n)\urcorner,
\]
such that $\alpha,\beta$ define $D_0,D_1$ in $\N$,
$\PA\vdash\forall x\,\neg(\alpha(x)\land\beta(x))$, and
\[
 n\in D_0\Rightarrow\PA\vdash\alpha(\bar n),\qquad
 n\in D_1\Rightarrow\PA\vdash\neg\alpha(\bar n).
\]
Consequently, for every c.e. $A\supseteq\PA$, the target reductions
\[
 F_{\N}(a_0)\le(\widehat F_{\N}\theta)^*\widehat F_{\N}(t_A),\qquad
 F_{\N}(a_1)\le(\widehat F_{\N}\theta)^*
                  (\widehat F_{\N}\nu)^*\widehat F_{\N}(t_A)
\]
hold, where $\widehat F_{\N}$ is the chosen extension to the
restriction splitting. They also hold for $F_M,\widehat F_M$ in
every $M\equiv\N$.
\end{lemma}
\begin{proof}
Use finite computation traces ending with the two output tags.
Determinism proves disjointness in PA. A true trace can be verified
by a finite PA proof; for the second tag, combine that proof with
the disjointness theorem. Numeral construction and substitution are
the morphisms realized in \S5. Extension includes these PA theorems
among the $A$-theorems. The resulting uniform arithmetic implications
are true in $\N$, so they hold in $M\equiv\N$.
\end{proof}
\source{Retain the finite-computation argument of v19 \S8, but state the inequalities in the target. The proof does not turn separate numeral proofs into a $\PA$-provable universal implication. A source version would require an additional uniform proof transformation; it is not needed for this plan.}

Let $\theta_{\mathrm{num}},\nu_{\mathrm{num}}:\N\to\N$ be the
chosen primitive recursive code functions realizing $\theta$ and
$\nu$ on their specified carriers. Thus
$\theta_{\mathrm{num}}(n)=\ulcorner\alpha(\bar n)\urcorner$,
and $\nu_{\mathrm{num}}$ forms the negation code. Inside arithmetic
formulas, these function symbols abbreviate their PA-provably
functional and total graph definitions.

\begin{lemma}[A partial proof-search lift]\label{lem:proof-search}
Let $A$ be a c.e. axiom presentation. There is a partial morphism
$\ell_A:\omega\rightharpoonup\omega+\omega$ in $S_{\PA}$
searching for a proof of the sentence coded by
$\theta_{\mathrm{num}}(x)$ or of its negation and returning
the reversed tag. In any model $M\models\PA$ in which these two
proof predicates are disjoint, its interpretation is the partial
conditional for their two domains.
\end{lemma}
\begin{proof}
Let $C(x)=\Prov_A(\theta_{\mathrm{num}}(x))$ and
$C'(x)=\Prov_A(\nu_{\mathrm{num}}(\theta_{\mathrm{num}}(x)))$.
Enumerate witnesses to the two
$\Sigma_1$ predicates in increasing finite bounds, with a fixed
priority if both are first found at the same stage. PA proves that
the search halts exactly on $C\lor C'$ and that every returned
value satisfies
\[
 (C(x)\land z=\iota_1(x))\lor
 (C'(x)\land z=\iota_0(x)).
\]
Here $\iota_j(x)=2x+j$ realizes the injection $i_j$.
In a model where $C,C'$ are disjoint, this relation is functional,
so the interpreted search has exactly the displayed graph.
Neither disjointness nor totality needs to be proved in PA.
\end{proof}
\source{Generalize v19 Lemma 7.2. Do not invoke the functional-search equality conclusion before functionality is available: the algorithm is defined first; exact equality with the disjunctive graph is proved in the target.}

\begin{theorem}[Incompleteness from the internal theory object]\label{thm:arithmetic-incomplete}
For every consistent c.e. extension $A\supseteq\PA$, the theory
object constructed in $\B_{\PA}$ has an incomplete standard
interpretation. There is a total source morphism
$g:\omega\to\omega$ whose standard interpretation
$\gamma=F_{\N}(g):\N\to\N$ satisfies: no number $\gamma(n)$ belongs to either
of the inverse images of its theorem and refutation sets under
the numerical substitution function $\theta_{\mathrm{num}}$. In particular,
$\alpha(\overline{\gamma(0)})$ is undecidable in $A$.
\end{theorem}
\begin{proof}
Use actual consistency to obtain disjoint target domains, or use
Corollary~\ref{cor:consistency} for an ultrapower interpretation.
Lemma~\ref{lem:reduction} supplies the target inequalities and
Lemma~\ref{lem:proof-search} supplies the partial lift.
Theorem~\ref{thm:logical} gives the total $g$ and the two zero
inverse images. For the individual ordinary sentence, use the
standard interpretation and evaluate at $0$.
\end{proof}
\source{New categorical derivation of the existing arithmetic conclusion through $S_{\PA}$. No completeness assumption is needed to construct an undecidable instance. The chosen $g$ may depend on the axiom presentation $A$.}

\begin{theorem}[Recursive inseparability and essential incompleteness]\label{thm:arithmetic-insep}
For every consistent c.e. extension $A\supseteq\PA$, its theorem
and refutation sets of sentence codes are recursively inseparable.
Every consistent c.e. extension of $A$ is therefore incomplete.
\end{theorem}
\begin{proof}
Suppose a recursive set separates the theorem and refutation sets.
Pull it back along the recursive numeral-substitution function
$\theta_{\mathrm{num}}$.
Its two complementary domains contain $D_0,D_1$ by
Lemma~\ref{lem:reduction}. Its reversed conditional is a total
recursive function. Choose a program index for it; that index is
a partial morphism of $S_{\PA}$, regardless of whether PA proves
its totality. Its standard interpretation is exactly the target
conditional. Corollary~\ref{cor:separator} gives a contradiction.
The same proof applies to every consistent c.e. extension. A
complete consistent c.e. theory would have complementary c.e.
theorem and refutation sets, whose simultaneous enumeration decides
them, contradicting recursive inseparability.
\end{proof}
\source{Recover v19 Corollary 8.2 and its recursive-inseparability conclusion. The full target category $\Pfn$ has arbitrary set-theoretic separators; only recursive separators supply the required program lifts.}

\paragraph{Extension maps.}
For c.e. axiom presentations $A\subseteq A'$ within the fixed
ambient $\B_{\PA}$, the generated-theory construction gives
$\Th_A\le\Th_{A'}$ whenever the chosen axiom inclusion is a
morphism of the source. A presentation of $A'$ containing the chosen
presentation of $A$ as an explicit tagged component gives a
PA-verifiable inclusion. A common PA component verifies only the
inclusions from PA into the two presentations. For the
external essential-incompleteness statement, ordinary inclusion of
theories and target inclusion suffice. The identity on the standard sentence
carrier gives a covariant map between the corresponding disjoint
theorem/refutation pairs. Do not replace this extension of represented
theories by the distinct quotient functor $S_B\to S_{B'}$.

\paragraph{Comparison with the numerical proof.}
Retain the $\mathrm{Rec}$ realization and explain that the standard
interpretation of the internal objects in $\B_{\PA}$ gives exactly
the same c.e. proof and theorem carriers. The new argument uses the
source Turing morphism and the interpreted source proof-search
morphism. It no longer leaves the syntactic-category construction
disconnected from the arithmetic incompleteness conclusion.

\section{Miller's sequence and a theory of real closed fields}
\purpose{Retain the application that distinguishes uniform computation from computation relative to a fixed representation.}
This section retains v19 results 4.1--4.6. It can follow the
arithmetic application because the comparison theorem and the theory
construction are then available. Its substantive conclusions do not
depend on the arithmetic ultrapower construction.

Fix a recursive pairing $\langle-,-\rangle:\N\times\N\to\N$
and an acceptable oracle-program enumeration $(\Phi_e^p)_e$.
Put $Q=[0,1]^{\N}$ and define the binary representation map
\[
 \rho:\{0,1\}^{\N}\to Q,\qquad
 \rho(p)(n)=\sum_{j\ge0}p(\langle n,j\rangle)2^{-j-1}.
\]
Write $P_z=\rho^{-1}(\{z\})$. Each $P_z$ is nonempty and compact.
Fix a recursive rational enumeration $(q_j)_j$. A program $d$
computes a real $r$ from $p$ when $\Phi_d^p(k)$ is defined for
every $k$ and $|q_{\Phi_d^p(k)}-r|<2^{-k}$. Choose a Miller
sequence $a\in Q$ with the index property: if $d$ computes the
same real $r\in[0,1]$ from every $p\in P_a$, then $a(d)=r$.
The index property is the strengthened conclusion following
Miller's Theorem~6.3; the binary-representation formulation is
given by Bauer--Hanson
\cite[Theorem~6.3 and following paragraph]{Miller}
\cite[\S2.1 and Theorem~2.2]{BauerHanson}.

\begin{lemma}[Unanimous computation]\label{lem:miller-unanimity}
There is one oracle program which, from $p\in P_z$ and numerical
inputs $e,n$, computes the partial function
$\tau_z:\N\times\N\rightharpoonup\N$ defined by
\[
 \tau_z(e,n)=m\quad\Longleftrightarrow\quad
 (\forall q\in P_z)\;\Phi_e^q(n)\mathord\downarrow=m.
\]
Its value and domain depend only on $z$, not on the selected
representation $p$.
\end{lemma}
\begin{proof}
Finite halting computations enumerate the open set
$H_{e,n,m}=\{q:\Phi_e^q(n)\mathord\downarrow=m\}$ by cylinders.
From $p\in P_z$, enumerate cylinders $[\sigma]$ for which
$z\notin\rho([\sigma])$: the image is a finite union of closed
dyadic boxes, and strict inequalities are semidecidable from $p$.
These cylinders enumerate the complement of $P_z$. For each $m$,
search for a finite cylinder cover of Cantor space using this
complement and $H_{e,n,m}$. The finite-cover test is decidable;
compactness makes success equivalent to $P_z\subseteq H_{e,n,m}$.
Dovetail these searches over $m$. Nonemptiness of $P_z$ ensures
uniqueness of the returned value, and the success condition depends
only on $z$.
\end{proof}

Let $\E_a$ have the objects generated from $\N$ by binary
products and tagged coproducts, with fixed recursive codings.
A morphism is a numerical partial function computed by one oracle
program with the same domain and values for every $p\in P_a$;
morphism equality is equality of partial functions.

\begin{proposition}[The uniform recursion category]\label{prop:miller-category}
The category $\E_a$ is a locally connected recursion category with
ranges, binary restriction coproducts, finite compatible joins,
and weak choice. The morphism
$\tau_a:\N\times\N\rightharpoonup\N$ is a Turing morphism.
\end{proposition}
\begin{proof}
Retain v19 Proposition 4.2: structural operations preserve uniform
domains and values; search gives ranges and sections; the preceding
compactness construction provides the evaluator and the parameter
theorem supplies its Turing factorizations. Singleton restrictions
detect equality for local connectedness.
\end{proof}

Choose the recursive coproduct isomorphism that tags a natural
number by its parity. The resulting diagonal domains are
\[
 A_i=\{e:\tau_a(e,e)\mathord\downarrow=m
                 \text{ for some }m\equiv i\pmod 2\},
 \qquad i=0,1.
\]
Let
$L_a=\{n:a(n)<1/2\}$ and $R_a=\{n:a(n)>1/2\}$.

\begin{lemma}[Reduction to the cut-index sets]\label{lem:miller-cuts}
There is an ordinary total recursive function $h:\N\to\N$ such
that $h^{-1}(L_a)=A_0$ and $h^{-1}(R_a)=A_1$.
\end{lemma}
\begin{proof}
Retain v19 Lemma 4.3: disjoint effectively open halting sets give
computable continuous functions positive exactly on those sets.
Their signed difference about $1/2$ yields a uniformly computable
real-valued function. Miller's index property gives the index $h(e)$
realizing its value at $a$ \cite{Miller}.
\end{proof}

\begin{corollary}[Inseparability and fixed-representation separation]\label{cor:miller-cuts}
The partial identities on $L_a,R_a$ are inseparable in $\E_a$.
For each fixed binary representation $p$ of $a$, a $p$-decidable
set separates the two underlying numerical sets.
\end{corollary}
\begin{proof}
Use source conditionals and Corollary~\ref{cor:source} for the
diagonal pair, then pull a hypothetical cut separator back along
$h$. The first binary digit of the fixed representation gives the
$p$-decidable separator. Thus the faithful range functor
$J_p:\E_a\to\mathrm{Rec}^{p}$ sends this pair to a separable
pair. If it preserved the chosen Turing morphism, the target
diagonal theorem would contradict that separator. Hence it does not.
\end{proof}

\begin{proposition}[First-order realization in the total splitting]\label{prop:miller-realization}
The category $\Tot(\Split\E_a)$ realizes the bounded-arity
first-order syntax, substitution, lists, and generated-theory
construction of Appendix A.
\end{proposition}
\begin{proof}
Retain v19 Proposition 4.5. Carriers are uniformly semidecidable
from all representations. Pairing, tags, equality, and uniform
search give the total categorical structure. Evaluate finite syntax
and derivation trees with the common operations; their domains and
values remain independent of the representation. Induction proves
initiality, and ranges give the generated theories.
\end{proof}

\begin{theorem}[Essential incompleteness in the uniform computability class]\label{thm:miller-incomplete}
Let $T_a$ consist of a recursive axiomatization of real closed
ordered fields in the language with constants $(c_n)_{n\in\N}$,
together with the axioms
\[
 q<c_n\quad(q<a(n)),\qquad c_n<q\quad(a(n)<q),
 \qquad q\in\mathbb Q, n\in\N.
\]
Every consistent extension of $T_a$ whose axiom set has a partial
identity belonging to $\E_a$ has inseparable theorem and refutation
domains in $\Split\E_a$, and is incomplete.
\end{theorem}
\begin{proof}
The real field with $c_n=a(n)$ gives consistency of $T_a$.
The realized proof object gives the theorem domains for each
extension. The map $n\mapsto\ulcorner c_n<1/2\urcorner$ reduces
$L_a,R_a$ to the theorem and refutation domains. Actual consistency
is source disjointness in this extensional category. Apply
Corollary~\ref{cor:source-logical}, retaining v19 Theorem 4.6.
\end{proof}
\source{Preserve the exact oracle representation hypotheses and the full cut-reduction proof from v19. No assertion that arbitrary target separators lift is made.}

\section{Three-point cosheaves and propositional theories}
\purpose{Retain a recursion category not restriction equivalent to $\mathrm{Rec}$ and explain the separate role of its two endpoint interpretations.}
Let $\mathcal C_{\mathrm{rec}}$ be the category of c.e. numerical
carriers and partial recursive functions total on their source
carriers, as in Proposition~\ref{prop:rec-realization}.
Use the fork diagram category $\E=\mathcal C_{\mathrm{rec}}^{J}$,
where $J$ has objects $0,1,2$ and nonidentity arrows $0\to1,0\to2$,
its partial-map category
$\A_X=\Par(\E,\operatorname{Mono})$, and the recursion subcategory
$\R_X$ constructed in Appendix B. Its splitting is equivalent
to $\A_X$. Identify $\E$ with $\Tot(\A_X)$ through the canonical
equivalence. The chosen object
$V\in\E$ has components and structure maps
\[
 \begin{gathered}
 V_0=\N\times(\N\times(\N\times\N)),\qquad
 V_1=V_2=\N\times\N,\\
 r_1^V(k,z,u,v)=(k,u),\qquad
 r_2^V(k,z,u,v)=(k,v).
 \end{gathered}
\]
Appendix B constructs a Turing morphism
$\tau:V\times V\rightharpoonup V$ and an isomorphism
$\zeta:V\to V+V$ whose component tags are determined by the
parity of the first coordinate. Put
\[
 a=(\zeta\tau\Delta_V)^*(1_V+0_V),\qquad
 b=(\zeta\tau\Delta_V)^*(0_V+1_V).
\]
Write $D_a,D_b\hookrightarrow V$ for their c.e. domain subfunctors
and put $A_j=(D_a)_j$, $B_j=(D_b)_j$, for $j=0,1,2$.
The full inclusion $\R_X\hookrightarrow\A_X$ lifts the
conditionals supplied by Proposition~\ref{prop:conditionals}.
Corollary~\ref{cor:source} makes $a,b$ inseparable; fullness retains
this conclusion in $\A_X$.
The object $V$ does not have decidable equality, so this application
uses the propositional construction of Appendix A.

Put $\mathcal K=\Par(\mathcal C_{\mathrm{rec}},\operatorname{Mono})$.
Component evaluation gives functors
$E_j:\A_X\to\mathcal K$ for $j=0,1,2$; on a span it takes
the span of its $j$th components.

\begin{proposition}[Endpoint interpretation]\label{prop:fork-endpoints}
The functors $E_j$ preserve restrictions, zeros, ranges, binary
restriction products and coproducts, and finite joins of restriction
idempotents. For $i=1,2$, the morphism
\[
 E_i(\tau):(\N\times\N)\times(\N\times\N)
                  \rightharpoonup\N\times\N
\]
is a Turing morphism in $\mathcal K$, and $A_i,B_i$ are recursively
inseparable. The morphism $E_0(\tau)$ is not Turing in $\mathcal K$.
\end{proposition}
\begin{proof}
Retain the explicit endpoint evaluator computations of v19
Proposition 9.1. Let $\mathcal K_2$ be the full subcategory of
$\mathcal K$ generated from $\N\times\N$ by the chosen binary
products and coproducts. Its objects are recursively isomorphic to
$\N\times\N$, and the displayed evaluator and its parameter maps
make it a recursion category. The underlying-partial-function
functor $\mathcal K_2\to\Pfn$ is a range functor.
Complementary c.e. domains have total recursive conditionals, so
Corollary~\ref{cor:source} gives the inseparability of the endpoint
diagonal pair. A recursive separator would be a complementary
restriction pair in $\mathcal K_2$.
For the zeroth component, retain the v19 counterexample
$f:V_0\to V_0$, $f(k,z,u,v)=(0,0,z,0)$: naturality would make
$r_1^Vf(k,z,u,v)$
independent of $z$ if $f$ were a specialization of $E_0(\tau)$.
\end{proof}

Let $S\in\E$ have finite propositional formula trees as its
components, with atom labels in $V_j$ and structure maps induced
by relabeling along $r_i^V$. Its constructors are morphisms
\[
 \theta:V\to S,\qquad \nu:S\to S,
 \qquad {\Rightarrow}:S\times S\to S.
\]
Use the three classical propositional Hilbert schemes specified in
v19 \S9, the positive and negative atomic axiom morphism
$[\theta|_{D_a},\nu\theta|_{D_b}]:D_a+D_b\to S$, and modus ponens.
The finite derivation trees give an object $\mathrm{Prf}\in\E$
and a conclusion morphism $c:\mathrm{Prf}\to S$.

\begin{proposition}[Propositional syntax and derivations]\label{prop:fork-initiality}
The constructor algebra of $S$ is initial for the endofunctor
$Z\mapsto V+Z+(Z\times Z)$ on $\E$. The derivation algebra
$(\mathrm{Prf},c)$ is initial for the rule endofunctor on $\E/S$
determined by these axioms and modus ponens.
\end{proposition}
\begin{proof}
Evaluate finite trees at each component. Induction proves that
evaluation commutes with the two fork structure maps, hence is a
morphism of $\E$; componentwise uniqueness proves initiality.
The same argument applies to derivations with their conclusion
labels. Retain v19 Propositions 9.2 and 9.3 as separate manuscript
statements.
\end{proof}

View $c$ as a total morphism in $\A_X$ and put $t_0=\ran(c)$,
a restriction idempotent on $S$. Let $D_{t_0}\hookrightarrow S$
be its domain subfunctor.

\begin{theorem}[Generated component theories]\label{thm:fork-generated}
The subfunctor $D_{t_0}$ is the least c.e. subfunctor of $S$
containing the logical and specified positive and negative atomic
axioms and closed under modus ponens. Its component theories are
consistent, and
\[
 a\le\theta^*t_0,\qquad b\le\theta^*\nu^*t_0,
 \qquad t_0\nu^*t_0=0_S.
\]
\end{theorem}
\begin{proof}
Apply Appendix A's least-theory and range/image results to the
initial derivation algebra. One-node proofs give the two
inequalities. For any alleged pair of contradictory finite proofs,
make precisely their positive atomic labels true. Disjointness of
the positive and negative atomic axioms makes every leaf valid;
induction contradicts opposite conclusions. This gives componentwise
consistency and hence source disjointness in $\A_X$.
\end{proof}

\begin{theorem}[Compatible extensions and endpoint incompleteness]\label{thm:fork-extensions}
Every compatible family of consistent c.e. deductively closed
extensions of the component theories has inseparable theorem and
refutation subobjects in $\A_X$. At each endpoint, their inverse
images under the atom map are disjoint recursively inseparable c.e.
sets. Extension inclusion gives the covariant maps of pairs
specified in v19 \S9, and each endpoint theory is incomplete.
\end{theorem}
\begin{proof}
Use Corollary~\ref{cor:source-logical} in $\A_X$ for the first
assertion. At endpoint $i$, the two c.e. inverse images contain
$A_i,B_i$. A recursive separator would separate that endpoint
diagonal pair. Completeness would make the inverse images
complementary c.e. sets, hence recursive. Monotonicity of both
sets under extension gives the asserted identity-on-carrier
morphisms of pairs.
\end{proof}

\begin{corollary}[Essential incompleteness of the endpoint theories]\label{cor:fork-essential}
Every consistent c.e. extension of either endpoint base theory is
incomplete.
\end{corollary}
\begin{proof}
Retain v19 Corollary 9.6 and its direct endpoint argument using
the same atomic reductions. This conclusion does not require every
endpoint extension to have been chosen as part of a compatible
family on the entire fork.
\end{proof}
\source{Retain all substantive claims of v19 \S9. The source theory and pair constructions are categorical; the endpoint conclusion uses the separately verified numerical interpretation.}

\appendix
\section{Syntax and deductively closed theories}
\purpose{Define the internal syntax and generated theories categorically. This appendix contains no incompleteness or numerical realization argument.}
Move v19 Appendix B here without cutting its dependencies. Its
six subsections retain the following order and functions.
\begin{enumerate}
\item \emph{Categorical hypotheses and notation.} Regularity,
  disjoint stable finite coproducts where needed, explicitly
  parenthesized finite powers, $\mathbf2=1+1$, Boolean morphisms
  defined by copairings, and conditional decompositions.
\item \emph{First-order syntax.} The ambient category $\D$, the
  variable object $V\in\D$ with decidable equality, bounded-arity
  signature objects, the paired syntax endofunctor
  $K(X,Y)=(K_1(X,Y),K_2(X,Y))$, parameterized initiality, and lists.
\item \emph{Substitution.} Free-variable and free-for constructions,
  term and formula substitution, quantifier cases, and their
  universal-property proofs.
\item \emph{Inference rules and closed subobjects.} Rule objects,
  premise and conclusion maps, the slice endofunctor, its initial
  derivation algebra, and the image argument.
\item \emph{First-order theories.} The specified logical axiom
  schemes, generalization, modus ponens, sentence pullback, and
  generated-theory closure operator.
\item \emph{Ranges and regular images; propositional instance.}
  The precise comparison with partial-map ranges and the
  propositional specialization needed by \S9.
\end{enumerate}
Retain three separate definitions from v19: Definition B.1,
\emph{categorical syntax}; Definition B.2, \emph{the sentence object};
and Definition B.4, \emph{a deductively closed subobject}.
The sentence-object definition follows the free-variable construction,
and the definition of deductive closure precedes the least-image
theorem. Together with the four results below, these account for
all seven numbered items in the current syntax appendix.

\begin{proposition}[Categorical substitution]
Under the appendix's parameterized syntax and list initiality and
decidable-variable hypotheses, the specified free-variable,
free-for, and substitution operations exist as morphisms of $\D$
and satisfy their constructor equations, including quantifiers.
\end{proposition}
\begin{proof}
Retain the enlarged-algebra construction used to keep the original
syntax component together with the recursively computed output.
Parameterized initiality gives the fold; its projection to syntax
is the identity by uniqueness. Conditional copairings use the
declared decidable predicates. This is v19 Proposition B.3.
\end{proof}

\begin{theorem}[Least rule-closed subobject]
Let $\D$ be regular with finite coproducts, let $\Phi\in\D$,
and let finitely bounded rule arities be given by
$p_n:R_n\to\Phi^n$ and $q_n:R_n\to\Phi$.
If the induced endofunctor $H:\D/\Phi\to\D/\Phi$ has an
initial algebra with conclusion $c:P\to\Phi$, then the regular
image of $c$ is the least subobject closed under these rules.
\end{theorem}
\begin{proof}
The rule endofunctor preserves regular epimorphisms: its summands
are pullbacks of finite powers, followed by a finite coproduct.
Factor $c=me$. The regular epimorphism $H(e)$ and the monomorphism
$m$ give a unique diagonal defining an algebra structure on the
image. Initiality maps $P$ into every other closed subobject;
the image factorization gives the required inclusion. Retain
v19 Theorem B.5 and its typed commutative square.
\end{proof}

\begin{theorem}[Generated first-order theory]
With the specified categorical first-order syntax and an axiom
subobject of the sentence object, an initial derivation algebra
gives the least formula subobject containing all axioms and closed
under the rules. Its sentence pullback is the least generated
sentence theory. When these derivation algebras exist for every
axiom subobject, the construction is an extensive, monotone,
idempotent closure operator on sentence subobjects.
\end{theorem}
\begin{proof}
Apply the preceding theorem to the specified rule objects.
Regular images commute with the sentence pullback. Leastness gives
extensivity and monotonicity; the generated theory is already
closed, which gives idempotence. Retain v19 Theorem B.6 and its
following closure argument.
\end{proof}

\begin{lemma}[Regular images and ranges]
In the partial-map presentation under the appendix's stated
admissibility hypotheses, the restriction idempotent splitting
through the regular image of a total map is its range. The
corresponding assertion for a partial map applies to its total
part on its domain.
\end{lemma}
\begin{proof}
The image idempotent receives the map. Every receiving restriction
idempotent supplies a monomorphism through which the total part
factors, so regular-image leastness makes the image idempotent
smaller. This is the range universal property. Retain the explicit
hypotheses and range-category references of v19 Lemma B.7
\cite{RangeII}.
\end{proof}
\source{Retain the existing literature dependencies on inductive definitions, initial-algebra syntax, binding, lists, and range categories where actually used. No numerical verification is moved into this appendix.}

\section{The three-point cosheaf recursion category}
\purpose{Supply the complete categorical construction used by \S9.}
Move v19 Appendix A here. Preserve the following ten statements
and their proofs; the changes are numbering and cross-references.
\begin{longtable}{@{}p{.13\textwidth}p{.30\textwidth}p{.50\textwidth}@{}}
\toprule
v19 & Statement to retain & Proof dependency and function\\\midrule\endhead
A.1 & Equivalence between cosheaves on the specified three-point space and fork diagrams. & Compute values on the opens and the cosheaf gluing conditions; give both functors.\\
A.2 & The category of fork diagrams of c.e. numerical carriers has the claimed finite-limit, coproduct, and image structure. & Componentwise constructions; verify the maps are recursive and satisfy the fork equations.\\
A.3 & The associated partial-map category has ranges, weak choice, and the stated restriction structure. & Use admissible monos and image factorizations from A.2; construct partial inverse sections.\\
A.4 & Local connectedness of that partial-map category. & Construct connected restrictions and prove that they detect every component of a partial morphism.\\
A.5 & Every object admits a restriction embedding into the chosen universal object. & Retain the component maps and the equations for their partial inverses.\\
A.6 & The coordinate frames and their compatibility equations. & Preserve the full data needed by the evaluator construction.\\
A.7 & The chosen evaluator is a Turing morphism. & Give the typed parameter morphism and verify all three component equations.\\
A.8 & The designated full subcategory is a recursion category whose restriction splitting recovers the partial-map category. & Combine A.3 and A.5--A.7 with isotypicality; use the restriction embeddings from A.5 for the splitting equivalence.\\
A.9 & Its Karoubi envelope also recovers the partial-map category. & Split arbitrary idempotents through their images and use the equivalence from A.8.\\
A.10 & It is not restriction equivalent to $\mathrm{Rec}$; the chosen atom object does not have decidable equality. & Retain the terminal-subobject count and the separate obstruction to decidable equality.\\
\bottomrule
\end{longtable}
\source{The exact definitions and proofs remain those of v19. The outline does not replace the component calculations needed by the recursion-category and non-equivalence claims.}

\section{Why lifting cannot be omitted}
\begin{proposition}[A range comparison without the required lifts]
There is a recursion-category example with a faithful range functor
preserving restriction products and reflecting zero maps for which
the source-inseparability conclusion fails when the section-lifting
hypothesis is removed.
\end{proposition}
\begin{proof}
Retain v19 Proposition C.1: use the category with oracle-enumerable
domains and ordinary partial recursive extensions, verify its
structural operations and evaluator, and exhibit the relevant
separators together with the missing lift. The new text should
identify Corollary~\ref{cor:source}, rather than an unconditional
claim about all target separators, as the statement whose omitted
hypothesis the example tests.
\end{proof}

\section{Faithful families of models and ambient theory extension}
\purpose{Retain the faithful representation theorem separately from the chosen soundness interpretation.}

\begin{theorem}[Faithful family of graph interpretations]
For consistent $B\supseteq\PA$, let $\mathcal M_B$ be the set
of arithmetic models of $B$ with underlying set $\N$, and put
\[
 \D_B=\prod_{M\in\mathcal M_B}\Pfn,
 \qquad F_B=(F_M)_M:S_B\to\D_B.
\]
Then $F_B$ is faithful and preserves the range, restriction,
product, coproduct, and finite-join structure stated in v19
Theorem 7.1. The target is split and locally connected, satisfies
weak choice and inverse-image preservation of joins, and
$F_B(1_\omega)\ne0$.
\end{theorem}
\begin{proof}
Preservation is coordinatewise. A graph equality not provable in
$B$ fails in a model, and completeness with downward
L\"owenheim--Skolem supplies a countable such model, transported
to the carrier $\N$. Singletons at individual coordinates detect
equality and give local connectedness. Partial inverse injections
give weak choice. Consistency makes the model family nonempty.
\end{proof}

\begin{proposition}[Ambient theory extension]
For consistent $B\subseteq B'$, the quotient functor
$Q_{B,B'}:S_B\to S_{B'}$ and restriction of model families
$\pi_{B,B'}:\D_B\to\D_{B'}$ give the commuting square
\[
\begin{tikzcd}
 S_B\arrow[r,"Q_{B,B'}"]\arrow[d,"F_B"'] & S_{B'}\arrow[d,"F_{B'}"]\\
 \D_B\arrow[r,"\pi_{B,B'}"'] & \D_{B'}.
\end{tikzcd}
\]
The quotient functor preserves the defining operations of the
chosen diagonal pair and sends it to the inseparable pair over
$B'$. These assignments compose for successive theory extensions.
\end{proposition}
\begin{proof}
Provability is monotone, so the quotient on indices is well defined.
Each coordinate of either composite interprets the same graph.
The quotient is surjective on hom-sets; a preimage of any target
morphism supplies its Turing factorization. Apply
Proposition~\ref{prop:st-pair} over $B'$.
\end{proof}
\source{Retain v19 Theorem 7.1, the extension square, and the extension part of Corollary 7.3. The image evaluator is not Turing among all target partial functions. Faithfulness and the product over all models are not used for soundness in \S6 or arithmetic incompleteness in \S7.}

\section*{Revision order and proof obligations}
\addcontentsline{toc}{section}{Revision order and proof obligations}
The next revision should proceed in the following order.
\begin{enumerate}
\item Replace the general comparison statement and prove its partial
  conditional version in \S3. Recover the existing source theorem
  as a corollary. Verify all coproduct comparisons and inverse-image
  equations before changing applications.
\item State the logical application in \S4 with target consistency,
  typed sentence-proof pullbacks, and only target reduction
  inequalities. Retain the direct source-consistency criterion for
  its actual uses.
\item Revise \S5 notation to distinguish ambient $B$ and represented
  $A$. Preserve the realization proofs and their parameterized
  universal properties. Add the explicit $B=\PA$ instantiation.
\item Write \S6, including the one-model soundness factorization,
  ultrapower comparison, and the distinction between the code sort
  and the semantic sort. Prove the exact standard proof-image
  identification for the chosen calculus.
\item Rewrite the arithmetic application in \S7 using partial
  proof search and the standard interpretation of $S_{\PA}$.
  Prove recursive inseparability using arbitrary recursive
  separator programs, without assuming their PA-provable totality.
\item Move and reconnect the Miller and cosheaf sections. Preserve
  their different computability classes, their actual consistency
  arguments, and the endpoint Turing calculations.
\item Reorder the appendices, move the faithful model family,
  update every cross-reference, and remove duplicated proof text
  only after its dependencies have been relocated.
\item Rewrite the introduction and abstract around the final
  statements. Compile and inspect the complete manuscript, then
  publish a consistently numbered manuscript release and source
  archive. The present plan does not change version 19's files.
\end{enumerate}

\paragraph{Claims requiring explicit verification during revision.}
The categorical diagonal and the ultrapower-factorization proofs
are given above. Integration still requires the exact program
implementation of the first-witness search, its PA-provable domain
and output implications, the identification of the selected
derivation-tree presentation with the stated proof predicates,
and the matching of all formula/term coding conventions. These are
implementation and exposition obligations for the prescribed proofs,
not assumptions to suppress. The target reductions can be proved
from ordinary finite computation proofs; the plan does not require
an unproved uniform source reflection principle.

\paragraph{Dependency summary.}
Appendix A and the realization in \S5 provide the internal theory.
Range preservation in \S4 identifies its interpreted proof image.
Section 6 supplies target consistency from a model, or transfers an
assumed actual consistency. The partial search in \S7 supplies the
lift; the comparison theorem in \S3 supplies the diagonal
contradiction or avoidance morphism. The faithful family in
Appendix D is independent of this route.

\section*{Literature used by this plan}
\addcontentsline{toc}{section}{Literature used by this plan}
The references below identify dependencies and comparison points for
the proposed organization. The full bibliography supporting the
retained syntax and cosheaf constructions remains in the manuscript.
The opening discussion of Di Paola--Heller and Keisler's stated
ultraproduct results were inspected for this plan. Detailed priority
claims about all categorical inseparability formulations are not
established by this outline.
\begin{thebibliography}{99}\small
\bibitem{BauerHanson} A.~Bauer and J.~E. Hanson, \emph{The Countable Reals}, \href{https://arxiv.org/abs/2404.01256v3}{arXiv:2404.01256v3}, 2026.
\bibitem{Bernardi} C.~Bernardi, \emph{Precomplete equivalence relations in dominical categories}, Proc. Amer. Math. Soc. \textbf{98} (1986), 124--134.
\bibitem{DPH} R.~A. Di Paola and A.~Heller, \emph{Dominical categories: recursion theory without elements}, J. Symbolic Logic \textbf{52} (1987), 594--635. \href{https://doi.org/10.2307/2274352}{doi:10.2307/2274352}.
\bibitem{Montagna} F.~Montagna, \emph{``Pathologies'' in two syntactic categories of partial maps}, Notre Dame J. Formal Logic \textbf{30} (1989), 105--116.
\bibitem{DPM} R.~A. Di Paola and F.~Montagna, \emph{Some properties of the syntactic $p$-recursion categories generated by consistent, recursively enumerable extensions of Peano arithmetic}, J. Symbolic Logic \textbf{56} (1991), 643--660.
\bibitem{CockettLack} J.~R.~B. Cockett and S.~Lack, \emph{Restriction categories. I. Categories of partial maps}, Theoret. Comput. Sci. \textbf{270} (2002), 223--259.
\bibitem{CockettHofstra} J.~R.~B. Cockett and P.~J.~W. Hofstra, \emph{Introduction to Turing categories}, Ann. Pure Appl. Logic \textbf{156} (2008), 183--209.
\bibitem{RangeII} J.~R.~B. Cockett, X.~Guo, and P.~Hofstra, \emph{Range categories II: Towards regularity}, Theory Appl. Categ. \textbf{26} (2012), 453--500. \href{https://www.tac.mta.ca/tac/volumes/26/18/26-18abs.html}{Journal text}.
\bibitem{Miller} J.~S. Miller, \emph{Degrees of unsolvability of continuous functions}, J. Symbolic Logic \textbf{69} (2004), 555--584. \href{https://doi.org/10.2178/jsl/1082418543}{doi:10.2178/jsl/1082418543}.
\bibitem{Keisler} H.~J. Keisler, \emph{The ultraproduct construction}, author's survey dated 1 June 2010, especially Theorem 3.1, Corollary 3.3, and Propositions 4.1--4.2. \href{https://people.math.wisc.edu/~hkeisler/ultraproducts-web-final.pdf}{Author's text}.
\bibitem{Terwijn} S.~Terwijn, \emph{The complexity of completions in partial combinatory algebra}, Math. Structures Comput. Sci. \textbf{34} (2024), 455--466, especially \S7. \href{https://doi.org/10.1017/S0960129524000252}{doi:10.1017/S0960129524000252}.
\bibitem{KLR} D.~E. Knuth, T.~Larrabee, and P.~M. Roberts, \emph{Mathematical Writing}, Stanford report STAN-CS-89-1193, 1989. \href{https://www-cs-faculty.stanford.edu/~knuth/papers/cs1193.pdf}{Stanford report}.
\end{thebibliography}
\end{document}
